% G-provenance.tex: граница заимствования из будущего в этом отчёте (добавлено 25.09.2026).
% Внешний читатель попросил перечислить каждый факт вне таблицы 16, узнанный после фиксации-снимка.
\section{Происхождение: что узнано после снимка}
\label{app:provenance}

Этот отчёт описывает систему на \code{1.0.0-rc.7}, фиксация \code{4e74376}. Правился он и после
этой фиксации, и кое-что из сказанного узнано позже. Это приложение есть полный перечень, полученный
из \code{git diff 4e74376} по исходникам отчёта, всего вне приложения~\ref{app:corrections}, чего один снимок читателю не
сказал бы. Каждая запись относится к одному из шести классов.

\begin{description}
  \item[Перемерено, воспроизводимо на снимке.] Числа, снятые позже командой над фиксацией
    \code{4e74376}; извлечь эту фиксацию и выполнить команду значит получить то же число: 313
    проверок, двенадцать мутационных учений, 80 сценариев учений, 45 миграций, 124 маршрута, 46
    триггеров, 44 поля вердикта, 2\,189 тестовых функций, из них 1\,030 продуктовых, 107 из 112
    криптографических тестов, 244 теста пакета OpenID4VP, шесть проверок Афины и разбивка стражей
    31/24/7.
  \item[Состояние репозитория или реестра, прочитанное позже.] Из фиксации не воспроизводится,
    потому что лежит вне её: три метки; четыре пакета и что каждый обслуживает;
    \code{полярис-проверка} на rc.3; таблица девяноста дней, прочитанная 23 сентября (день десятый,
    пять из семи); два внешних исправления.
  \item[Измерено 23 сентября.] Числа исполнения, снятые после снимка на той же машине: 2\,080,
    7\,950 и 740 проверок в секунду в разделе о ёмкости и времена доказательства 25\,мс и 12\,мс.
  \item[События после снимка.] Сертификация OpenID 24 сентября (разделы о верификаторах и
    ограничениях, приложение~A, приложение~F, рисунок о внешних свидетелях); правило объявления
    стражей от 23 сентября; расширенный охват учения процедур и его тест блокировки; измерение
    удаления триггеров (42, 38, 42); учение отделённого верификатора (45 отказов, 32 выжили, семь
    объявлены); порядок смены паролей; путь постороннего и README пакетов, пройденные 23 сентября;
    и находка, что половина доказательства в измерении ёмкости ничего не измеряла с
    \code{в9.354}.
  \item[Исправления того, что было неверно уже на снимке.] 143 теста пакета суть 142; 43 поля
    вердикта суть 44; ``пять проверок'' суть шесть проверок Афины; ``несколько сотен сводок
    кластеров'' суть не больше пяти тысяч; подпись к жизненному циклу (из резерва в отозванный);
    подпись к районам (семь таблиц штрихом); формулировка о столбцах идентификаторов; дата второго
    внешнего исправления.
  \item[Написано 25 сентября, после внешнего рецензирования.] Приложение~\ref{sec:lineage} карты и список
    литературы; это приложение; таблица непроверенных утверждений (таблица~\ref{tab:untested});
    отзыв слов ``устойчивость к атаке Сивиллы'' в пользу уникальности удостоверения в пределах
    записанной личности; оговорка к слову ``проверены'' для национальных целей; различение
    свидетелей в разделе~\ref{sec:privacy}; минимальное множество анонимности и его поведение при
    отказе (изменение кода того же дня, в приложении~\ref{app:corrections}); пять процентов как параметр управления;
    владелец базы данных как несвязанный участник; и кто выбирает автономное окно.
\end{description}

Редакционные изменения (вёрстка, изображение на обложке, сокращения глоссария и перештамповка
\code{rc.1} на \code{rc.7}) не несут сведений и не перечислены. Числа кода рядом с \code{rc.7} в
тексте даны на снимке; числа в README и на сайте даны на текущем дереве и от них отличаются.
